\chapter{Summary}
The evolution of eukaryotes was shaped by the phenomenon of endosymbiosis, which gave rise to the canonical organelles, the mitochondria and plastids, billions of years ago. Notably, this process was accompanied by the establishment of mechanisms for protein targeting and import into these canonical organelles. Organisms that are in the early stage of organellogenesis offer, therefore, great insights into the principles guiding these mechanisms. 
In this work, nitroplast targeting signals (known as \glsxtrshort{utp}s) from the algae \glsxtrlong{Bbig} were investigated. The aim was to understand the influence of these \glsxtrshort{utp}s on other endosymbiotic systems and to find suitable expression conditions to purify \glsxtrshort{utp}s for structural analysis. For that purpose, different fusion proteins containing these transit peptides were generated and integrated into different model organisms. Additionally, small bioinformatic analyses were conducted to further derive the mechanism of how \glsxtrshort{Bbig} encoded proteins arrive in the nitroplast. In the scope of this thesis, only three \glsxtrshort{utp}s were examined (\glsxtrshort{utp_B}, \glsxtrshort{utp_E}, \glsxtrshort{utp_C}). 
First, \glsxtrlong{Adea} and its endosymbiont were the model of choice for studying the impacts of the nitroplast targeting signals inside a heterologous system. The recombinant \glsxtrshort{utp}s expressed in \glsxtrshort{Adea} contained an \glsxtrshort{egfp}-tag for localization and verification purposes. Heterologous expression of the \glsxtrshort{egfp}-\glsxtrshort{utp} constructs, although weak, was confirmed. No conclusive data were collected on the localization of the nitroplast transit peptides within the cells. Importantly, the \glsxtrshort{utp}s had no impact on \glsxtrshort{Adea} or its endosymbiont.   
Second, the \glsxtrshort{utp}s were also expressed in \glsxtrlong{Eco} with an affinity tag for purification via \glsxtrlong{IMAC}. Expression, solubility, and purification tests were conducted. \glsxtrshort{Eco} was able to express soluble \glsxtrshort{utp}s when it was grown at 28~°C for 4~hours after induction. Expression cultures grown overnight or at temperatures above and below 28~°C had decreased amounts of soluble fusion proteins, either due to degradation or the formation of inclusion bodies. Furthermore, the \glsxtrshort{utp}s were able to be purified using \glsxtrshort{IMAC}.
Lastly, bioinformatic predictions showed that all three nitroplast targeting signals shared a bipartite structure containing an N-terminal region with amphipathic α-helical structures arranged in a U-shape and a more variable, likely flexible C-terminal region containing a transmembrane helix. It was also inferred that the mechanism for the transport of proteins targeted to the nitroplast might be similar to the way it likely occurs for the chromatophore.  
Overall, the expression of the fusion proteins containing the nitroplast targeting signals was successful in all heterologous systems. 
Going forward, simultaneous expression of \glsxtrshort{egfp}-tagged \glsxtrshort{utp}s in \glsxtrshort{Adea} from multiple loci and the utilization of the immunofluorescence assay would likely increase the chances of localizing the constructs in the cells.
Further, it was concluded that larger-scale expression and purification of \glsxtrshort{6xhis}-\glsxtrshort{utp}s using \glsxtrlong{fplc} to obtain enough purified nitroplast transit peptides for structural analysis should be possible. 
And lastly, the search for structural homologs, to infer possible functions or interactions of the two domains, should be continued using Foldseek. Further analyses of the motifs within the  \glsxtrshort{utp}s would potentially uncover more of the mechanism behind the targeting and transport to the nitroplast